\begin{table}[h]
\centering
\begin{tabular}{ll}
\hline\hline
\noalign{\smallskip}
{\bf V0 selection criterion} & Value                   \\
\noalign{\smallskip}
\hline
\noalign{\smallskip}
DCA (h$^\pm$ to PV)                        & $>0.06$  cm                           \\ 
DCA (h$^-$ to h$^+$)                       & $<1.0$ standard deviations                  \\
Fiducial volume (R$_{2D}$)                 & $>$ 0.5 cm                    \\
V0 pointing angle                          & $\cos\theta_{V0}>$ 0.97 (0.995)                     \\
Proper Lifetime                         & $<$ 20 (30) cm/$c$                     \\
Competing V0 Rejection Window               &  $\pm$5 (10) \MeVmass                     \\
\noalign{\smallskip}
\hline\hline
\end{tabular}
~\newline
\caption{Selection criteria parameters utilized in the \kzero, \lmb\ and \almb\ analyses presented in
this work. If a criterion for \lmb\ and \kzero\ finding differs, the criterion for the \lmb\ hypothesis
is given in parentheses. The acronym 
DCA stands for ``distance of closest approach" and PV for ``primary event vertex". The pointing angle
$\theta$ is the angle between the momentum vector of the reconstructed V0, and the line
segment bound by the decay and primary vertices and R$_{2D}$ denotes the transverse distance from the 
detector center. 
\label{tab:V0Sels}
}
\end{table}
